\section{Discussion and conclusions}\label{sec:discussion}

\subsection{What was shown}\label{sec:discussion-shown}

Within the models used, the study establishes six points. (i) The passive SiN elements meet the proposal's
targets on the 400~nm film, and two targets need restating: the steering update rate should be $\geq 200$~kHz
for moves at 1.2~mK, and position noise must be specified as a spectral density at the trap frequency rather
than as an rms value (\cref{sec:platform}). (ii) Independent motion of about 20 tweezers along a row, the
integrated equivalent of the acousto-optic deflector, is possible with a sliding-window focal-plane OPA. In
its optimised form (design D1: 128 ports, 192 arms and emitters at 1.5~\textmu m pitch), a 1.2~mK tweezer needs
about 57~mW once the facet loss and the glued collimator are included (\cref{sec:engine,sec:rows}). (iii) Rows of
such engines can be stacked and relayed onto a 4~\textmu m lattice. The rows do not need one frequency shifter
per chip: common phase noise on the seed (about 3~rad rms, 1 to 2~GHz) with different fibre delays makes them
mutually incoherent. Its conversion to intensity noise stays 16~dB inside the heating criterion if no filter
narrower than about 100~pm and no pair of reflections with $R_1R_2>1.7\times10^{-3}$ lies in the path
(\cref{sec:rows}). (iv) Transport with handoff in three dimensions delivers 3,000 of 3,000 simulated atoms when
the moves run along lanes between rows, the moving depth is at least three times the hold depth, and the
hold depth is about $20\,k_BT$. With 200~\textmu s ramps and 400~\textmu s lane steps the heating per handoff
falls to 0.12~\textmu K (about 0.3~\textmu K per move including the $x$ move) at a moving depth of only 0.3~mK, and the axial registration of the two foci
must be held to $\pm0.5$~\textmu m (\cref{sec:transport}). (v) At $10^6$ sites, a block of 128 engines with fast
steps holds a time-averaged defect fraction of $1.8\times10^{-4}$ to $2.0\times10^{-4}$, set by vacuum and
imaging losses rather than by the repair. Slower steps reach a similar floor if the imaging period is raised
to match the sweep time of the block (\cref{sec:megascale}). (vi) A refractive cylindrical collimator serves only the
central half of a segment, but two mirrors in a 0.48~mm fused-silica block glued to the facet collimate exactly at
every $x$ angle and obey the sine condition. In 1D the peak stays at or above 0.98 out to the slit edge and the
crosstalk at 4~\textmu m at or below $1.0\times10^{-4}$, so one engine serves its whole 180~\textmu m segment
(\cref{sec:combine-reflective}).

\subsection{What is model-dependent}\label{sec:discussion-model}

Several results rest on models whose limits matter. The x response of the engine comes from a 2D beam
propagation calibrated to an effective index, and the facet transmission from a 2D FDTD. The y axis of the
relay is scalar, with ideal relay lenses downstream of the collimator. The transport model is classical in
three dimensions, checked quantum mechanically in one dimension per axis. The steady-state simulator draws
losses at assumed rates (vacuum lifetime, loss per image and per move) from a single random seed per case.

The most consequential model dependence was found late. Every y-axis result before \cref{sec:combine} assumed
light without an x component. A tweezer at $x_a$ is formed by light with x direction sine $s_x = x_a/x_\mathrm{scale}$.
The refractive cylindrical collimator, which conserves $k_x$, then acts with effective indices $\sqrt{n^2-s_x^2}$
and defocuses the facet by about $-99\,s_x^2$~\textmu m, so a row is usable over $|x_a|\leq 45$~\textmu m only. The
all-reflective collimator removes this dependence exactly. In its place it brings two dependences of its own: the
mirror form, which must be held to about 10~nm rms, and the s-p phase of the coatings, which makes the tweezer
polarisation partly circular away from the row centre ($|V|\leq0.04$ at 90~\textmu m with a silvered or gilded
first mirror and the second in total internal reflection, 0.22 with two silvered mirrors). The 2D checks reach
60~\textmu m, the limit of the D1 engine used in the 2D model.

In addition, a review of the whole project for errors corrected several earlier results, all reflected in this
report: (i) the steady-state simulator skipped block positions after each visit, which had inflated every
assembly time and the slow-step floors at $10^6$ sites; (ii) the phase-noise option for row decorrelation
scaled away the incoherent part of the crosstalk and looked better than fully incoherent rows; (iii) the
$2\times2$ MMI design offset of the switches was treated as random spread; (iv) the field-mean laser power of
the first engine was evaluated in the wrong units; (v) the arm-phase aberration of the first star coupler was
mis-evaluated by phase unwrapping; (vi) the first $x$ relay lens did not cover the slit edge; (vii) the
position error of the 20-port star-coupler engine depends on the relative port phases, which an earlier text
treated as fixed; and (viii) the 2D values beyond 67~\textmu m in the first skew study were read as a break-up of the
spot, but lie outside the reach of the model's engine. Some figures still differ between documents for stated reasons: the heating per slow step is
0.12~\textmu K per handoff and about 0.3~\textmu K per move including the $x$ move (the 3D classical run gives
0.16~\textmu K); the greedy scheduler's pass count depends on the set of random loadings (7.8 or 8.05 passes);
the steering half-range at 128 channels is 67~\textmu m by the usable-range criterion and 68~\textmu m by the
formula $N\lambda/(4\,\mathrm{NA})$; the other tweezers shift by 1.5~nm in the paraxial engine and up to 3.5~nm in
D1; and the time-averaged floor at $10^6$ sites with fast steps is 1.8 to $2.0\times10^{-4}$ across runs with
different seeds. Defect fractions quoted at the image instant count imaging losses twice and exceed the time
average by about 25\% at the nominal point.

\subsection{What is extrapolated}\label{sec:discussion-extrap}

The large-array chapters extrapolate in three respects. (i) The optics for $10^6$ sites assume a corrected field
of about $\pm2.8$~mm at NA 0.5, which we did not design. (ii) The static potential for $10^6$ sites needs of the
order of kilowatts of laser power, from tweezers or from a cavity lattice, which we priced but did not
design. (iii) The comb that combines four 16-chip stacks into a 64-row block exists only as a first-order concept with
wave-optics checks of its key penalties. The reflective collimator is constructed and traced exactly, but its
coatings are not designed, its tolerances are evaluated one at a time, and no fabrication route is
demonstrated.

\subsection{Conclusions}\label{sec:discussion-concl}

We have shown, in simulation, that SiN circuits can steer independent transport tweezers along rows of a
4~\textmu m lattice, that stacked rows can be made independent without per-chip frequency shifters, and that
handoff transport is clean and cold with slow steps and lanes between rows. A refractive cylindrical collimator
serves only the central half of each engine's segment; a two-mirror collimator in glass, designed here, serves the
whole segment. The main limitation is now that collimator's fabrication: two freeform mirrors held to about 10~nm
rms in a 0.48~mm block glued at $\pm0.1$~\textmu m, with coatings of small s-p retardance. The next steps are a
criterion for the circular polarisation of the transport tweezers, a coating design, and measurements on
fabricated D1 engines of the quantities the models assume: phase accuracy and update rate of the actuators,
facet transmission, and the intensity noise of the decorrelated laser chain.
